\begin{proposition}[Two projective torus fibrations]\hypertarget{tag-01G3}{}\label{cdp:st:D.1.5}
\StatementRecord{cdp:st:D.1.5}{proved}{01G3}{\citet[Section~10.5, p.~91]{AC2026}}
Let $E_0=F_0=E'=\C/(\Z+i\Z)$ and $a=1/2\in E_0$.
The involution
\[
(e,z)\longmapsto(e+a,-z)
\]
acts freely on $E_0\times F_0$. Put
$S=(E_0\times F_0)/\langle(a,-1)\rangle$,
$Y=S\times E'$, and $C=E_0/\langle a\rangle$.
The induced map $Y\to C$ is a smooth torus fibration with fibre
$F_0\times E'$ and monodromy
$\operatorname{diag}(-I_2,I_2)$ around the nontrivial deck character.
Its integral homology coinvariants are
$\Z^2\oplus(\Z/2\Z)^2$. After removing $s\geq1$ fibres,
\[
H_1(U;\Z)=\Z^{s+3}\oplus(\Z/2\Z)^2.
\]
The sum of the first Betti numbers of the punctured base and the
fibre is $s+5$.

Let $\eta$ be the nontrivial order-two degree-zero line bundle on
$C$ associated to the double cover $E_0\to C$. The coherent direct
images are
\[
R^1f_*\OO_Y=\eta\oplus\OO_C,\qquad R^2f_*\OO_Y=\eta.
\]
The cohomology dimensions are $h^q(Y,\OO_Y)=(1,2,1,0)$.
Restriction from $H^1(Y,\OO_Y)$ to the two-dimensional
$H^1(F_0\times E',\OO)$ has image precisely the $E'$ direction
and has rank one. These homomorphisms are induced by the finite
\'etale quotient and projection, and $b_1(Y)=4$.

For an example over $\PP^1$, let
$h:R\to\PP^1$ be the relatively minimal elliptic surface with
Weierstrass equation
\[
y^2=x^3+(t^4+1)x+t^6.
\]
Its fundamental line bundle is $\OO(1)$ and its discriminant has
twelve simple zeroes. Moreover, $R^1h_*\OO_R=\OO(-1)$.
For $Y'=R\times E'\to\PP^1$, the product K\"unneth maps give
\[
R^1f_*\OO_{Y'}=\OO\oplus\OO(-1),\qquad
R^2f_*\OO_{Y'}=\OO(-1),
\]
and restriction in degree one again has rank one. This projective
threefold has $b_1=2$, $b_2=11$ and $a=3$.
\end{proposition}
\begin{proof}
The first construction specializes the bielliptic example of
\citet[Section~10.5, p.~91]{AC2026} to the displayed involution.
Translation by $a$ has no fixed point, so the quotient is smooth.
The map $[e]\mapsto([e,0],0)$ is a section of the fibration. Over
the punctured base its fundamental group is therefore the semidirect
product of the fibre lattice and the free base group of rank $s+1$.
The map $(k,q)\mapsto([q],[k])$ identifies its abelianization with
the direct sum of the base abelianization and the fibre coinvariants:
the relations are $\rho(q)k-k$, which here are twice the two
$F_0$ generators. The same section over the complete elliptic base
gives the stated value of the first Betti number.

On the covering product, the three antiholomorphic one-form
directions have signs $(1,-1,1)$. Descent of the fibre cohomology
gives the two direct images, and taking invariants in the exterior
algebra gives the stated cohomology dimensions. The base direction
restricts to zero on a fibre, while the invariant fibre direction is
exactly that of $E'$. This proves the description of the restriction
map.

For the Weierstrass surface, the coordinates $u=1/t$,
$x'=u^2x$, $y'=u^3y$ give
\[
y'^2=x'^3+(1+u^4)x'+1.
\]
The relative residue differential transforms by
$dx'/(2y')=u^{-1}dx/(2y)$, so the fundamental line is
$\OO_{\PP^1}(1)$. Its discriminant is
\[
-16P(t),\qquad P(t)=31t^{12}+12t^8+12t^4+4=Q(t^4),
\qquad Q(z)=31z^3+12z^2+12z+4.
\]
The identity
\[
\left(\frac5{36}-\frac{31z}{54}\right)Q(z)
+\left(\frac1{27}-\frac{7z}{324}+\frac{31z^2}{162}\right)Q'(z)=1
\]
shows that the roots of $Q$ are simple; its constant term is
nonzero. Thus $P$ has twelve simple zeroes. At infinity the
transformed discriminant is $-16\cdot31$, so the fibre there is
smooth. The twelve singular fibres are nodal cubics. Relative
duality gives the stated first direct image. Leray gives
$h^1(R,\OO_R)=h^2(R,\OO_R)=0$; the twelve nodal fibres give
$e(R)=12$, and hence $b_1(R)=0$ and $b_2(R)=10$.
K\"unneth now gives both direct images and the Betti numbers of
the product. Its degree-one coherent cohomology is the $E'$
summand, giving the restriction map. The product is projective
of dimension three, so its algebraic dimension is three.
\end{proof}